\documentclass{amsart}
% For dealing with references we use the comment environment
\usepackage{verbatim}
\newenvironment{reference}{\comment}{\endcomment}
%\newenvironment{reference}{}{}
\newenvironment{slogan}{\comment}{\endcomment}
\newenvironment{history}{\comment}{\endcomment}

% For commutative diagrams we use Xy-pic
\usepackage[all]{xy}

% We use 2cell for 2-commutative diagrams.
\xyoption{2cell}
\UseAllTwocells

% We use multicol for the list of chapters between chapters
\usepackage{multicol}

% This is generally recommended for better output
\usepackage{lmodern}
\usepackage[T1]{fontenc}

% For cross-file-references

% Package for hypertext links:
\usepackage{hyperref}

% For any local file, say "hello.tex" you want to link to please

% Theorem environments.
%
\theoremstyle{plain}
\newtheorem{theorem}[subsection]{Theorem}
\newtheorem{proposition}[subsection]{Proposition}
\newtheorem{lemma}[subsection]{Lemma}

\theoremstyle{definition}
\newtheorem{definition}[subsection]{Definition}
\newtheorem{example}[subsection]{Example}
\newtheorem{exercise}[subsection]{Exercise}
\newtheorem{situation}[subsection]{Situation}

\theoremstyle{remark}
\newtheorem{remark}[subsection]{Remark}
\newtheorem{remarks}[subsection]{Remarks}

\numberwithin{equation}{subsection}

% Macros
%
\def\lim{\mathop{\mathrm{lim}}\nolimits}
\def\colim{\mathop{\mathrm{colim}}\nolimits}
\def\Spec{\mathop{\mathrm{Spec}}}
\def\Hom{\mathop{\mathrm{Hom}}\nolimits}
\def\Ext{\mathop{\mathrm{Ext}}\nolimits}
\def\SheafHom{\mathop{\mathcal{H}\!\mathit{om}}\nolimits}
\def\SheafExt{\mathop{\mathcal{E}\!\mathit{xt}}\nolimits}
\def\Sch{\mathit{Sch}}
\def\Mor{\mathop{\mathrm{Mor}}\nolimits}
\def\Ob{\mathop{\mathrm{Ob}}\nolimits}
\def\Sh{\mathop{\mathit{Sh}}\nolimits}
\def\NL{\mathop{N\!L}\nolimits}
\def\CH{\mathop{\mathrm{CH}}\nolimits}
\def\proetale{{pro\text{-}\acute{e}tale}}
\def\etale{{\acute{e}tale}}
\def\QCoh{\mathit{QCoh}}
\def\Ker{\mathop{\mathrm{Ker}}}
\def\Im{\mathop{\mathrm{Im}}}
\def\Coker{\mathop{\mathrm{Coker}}}
\def\Coim{\mathop{\mathrm{Coim}}}

% Boxtimes
%
\DeclareMathSymbol{\boxtimes}{\mathbin}{AMSa}{"02}

%
% Macros for moduli stacks/spaces
%
\def\QCohstack{\mathcal{QC}\!\mathit{oh}}
\def\Cohstack{\mathcal{C}\!\mathit{oh}}
\def\Spacesstack{\mathcal{S}\!\mathit{paces}}
\def\Quotfunctor{\mathrm{Quot}}
\def\Hilbfunctor{\mathrm{Hilb}}
\def\Curvesstack{\mathcal{C}\!\mathit{urves}}
\def\Polarizedstack{\mathcal{P}\!\mathit{olarized}}
\def\Complexesstack{\mathcal{C}\!\mathit{omplexes}}
% \Pic is the operator that assigns to X its picard group, usage \Pic(X)
% \Picardstack_{X/B} denotes the Picard stack of X over B
% \Picardfunctor_{X/B} denotes the Picard functor of X over B
\def\Pic{\mathop{\mathrm{Pic}}\nolimits}
\def\Picardstack{\mathcal{P}\!\mathit{ic}}
\def\Picardfunctor{\mathrm{Pic}}
\def\Deformationcategory{\mathcal{D}\!\mathit{ef}}

\setlength{\emergencystretch}{3em}
\begin{document}
\title[Stacks Project, Tag 0AW3]{418 --- Failure of formal catenarity obstructs universal catenarity}
\maketitle
\noindent Copyright (C) 2005--2025 Johan de Jong. Stacks Project authors. Source: \href{https://stacks.math.columbia.edu/tag/0AW3}{Tag 0AW3}.

\noindent Editorial difficulty note (not author text). Difficulty H2: The proof chooses a smallest-dimensional completed component and algebraizes its quotient to contradict the dimension formula. For larger component dimension it chooses a perturbing parameter by prime avoidance and descends the defect inductively..

\noindent Editorial provenance note (not author text). History: excerpt from revision \texttt{a0444\allowbreak 6e57e\allowbreak c1fbc\allowbreak 252a8\allowbreak 71afc\allowbreak ec775\allowbreak 2fb28\allowbreak 07b14}, more-algebra.tex, lines 32324--32386. Reference presentation was adapted; original-author context and core support blocks were retained or added. The mathematical wording is unchanged. Licensed under GNU FDL 1.2 or later; no invariant sections or cover texts. See ../licenses/COPYING. Context blocks reproduce author definitions and supporting results; they are not additional counted problems.

\begin{lemma}
\label{lemma-not-formally-catenary}
Let $(A, \mathfrak m)$ be a Noetherian local ring which is not
formally catenary. Then $A$ is not universally catenary.
\end{lemma}

\begin{proof}
By assumption there exists a minimal prime $\mathfrak p \subset A$ such that
the spectrum of $A^\wedge /\mathfrak p A^\wedge$ is not equidimensional.
After replacing $A$ by $A/\mathfrak p$ we may assume that $A$
is a domain and that the spectrum of $A^\wedge$ is not equidimensional.
Let $\mathfrak q$ be a minimal prime of
$A^\wedge$ such that $d = \dim(A^\wedge/\mathfrak q)$
is minimal and hence $0 < d < \dim(A)$. We prove the lemma by induction
on $d$.

\medskip\noindent
The case $d = 1$. In this case $\dim(A^\wedge_\mathfrak q) = 0$.
Hence $A^\wedge_\mathfrak q$ is Artinian local and we see that
for some $n > 0$ the ideal $J = \mathfrak q^n$ maps to zero in
$A^\wedge_\mathfrak q$. It follows that $\mathfrak m$ is the
only associated prime of $J/J^2$, whence $\mathfrak m^m$ annihilates
$J/J^2$ for some $m > 0$. Thus we can use
Lemma \href{https://stacks.math.columbia.edu/tag/0ALR}{lemma quotient by idempotent (Tag 0ALR)}
to find $A \to B$ of finite type such that $B^\wedge \cong A^\wedge/J$.
It follows that $\mathfrak m_B = \sqrt{\mathfrak mB}$ is a maximal
ideal with the same residue field as $\mathfrak m$ and $B^\wedge$
is the $\mathfrak m_B$-adic completion
(Algebra, Lemma \href{https://stacks.math.columbia.edu/tag/0394}{lemma finite after completion (Tag 0394)}).
Then
$$
\dim(B_{\mathfrak m_B}) = \dim(B^\wedge) = 1 = d.
$$
Since we have the factorization $A \to B \to A^\wedge/J$ the inverse image
of $\mathfrak q/J$ is a prime $\mathfrak q' \subset \mathfrak m_B$ lying
over $(0)$ in $A$. Thus, if $A$ were universally catenary, the dimension
formula (Algebra, Lemma \href{https://stacks.math.columbia.edu/tag/02IJ}{lemma dimension formula (Tag 02IJ)}) would give
\begin{align*}
\dim(B_{\mathfrak m_B})
& \geq
\dim((B/\mathfrak q')_{\mathfrak m_B}) \\
& =
\dim(A) + \text{trdeg}_A(B/\mathfrak q') -
\text{trdeg}_{\kappa(\mathfrak m)}(\kappa(\mathfrak m_B)) \\
& =
\dim(A) + \text{trdeg}_A(B/\mathfrak q')
\end{align*}
This contradiction finishes the argument in case $d = 1$.

\medskip\noindent
Assume $d > 1$. Let $Z \subset \Spec(A^\wedge)$ be the union of
the irreducible components distinct from $V(\mathfrak q)$.
Let $\mathfrak r_1, \ldots, \mathfrak r_m \subset A^\wedge$
be the prime ideals corresponding to irreducible components of
$V(\mathfrak q) \cap Z$ of dimension $> 0$.
Choose $f \in \mathfrak m$, $f \not \in A \cap \mathfrak r_j$
using prime avoidance (Algebra, Lemma \href{https://stacks.math.columbia.edu/tag/00DS}{lemma silly (Tag 00DS)}).
Then $\dim(A/fA) = \dim(A) - 1$ and there is some irreducible
component of $V(\mathfrak q, f)$ of dimension $d - 1$.
Thus $A/fA$ is not formally catenary and the invariant $d$ has
decreased. By induction $A/fA$ is not universally catenary, hence
$A$ is not universally catenary.
\end{proof}
\end{document}
